\section{An internal sale becomes group profit only through an outside transaction}
\label{sec:worked-intragroup}

The acquisition worksheet now has a second task: keep two valid legal books
from counting a transaction within the group as income from outside it. Continue
the synthetic 75 per cent cash acquisition, but add inventory costing 20 and
matching capital of 20 to the parent's opening balances. After obtaining
control, the parent sells the whole lot to the subsidiary for 30 on credit.
The parent records revenue 30, cost of sales 20 and a receivable 30. The
subsidiary records inventory 30 and a payable 30. Each legal book balances,
yet their sum includes a profit of 10 before any customer outside the group
has bought the goods.

Consolidation removes the reciprocal receivable and payable in full. It also
removes the internal revenue 30 against cost of sales 30, then charges the
unrealised margin 10 to cost of sales and reduces inventory by 10. After these
two profit adjustments, the parent's original cost of sales 20 contributes
zero, and the group carries inventory at its original cost of 20. Multiplying
the adjustments by the parent's 75 per cent ownership would leave internal
profit in group income. Full elimination and attribution of the subsidiary's
result are separate operations \citep[paras.~B86--B89]{aasb10control2026}.

The original acquisition adjustments still apply. Equipment retains the
acquired basis and the additional depreciation; the subsidiary's pre-control
profit remains outside current group profit. The internal sale takes place
after control, so it needs no retrospective reassignment of a pre-control
transaction. If the subsidiary subsequently pays the 30, its cash payment and
the parent's cash receipt cancel as separate intragroup flows. Settlement
removes the reciprocal balances from the legal books, but does not realise
the inventory profit for the group.

\subsection{Removing profit does not remove an external tax obligation}

Tax depends on more than the consolidated carrying amount. AASB 112 compares
that amount with the appropriate tax base, which may depend on consolidated
or separate tax returns. Its illustrative examples include intragroup profit
eliminated from inventory without an equivalent adjustment for tax purposes
\citep[para.~11 and illustrative examples, B11]{aasb112dated}. The accounting
rule does not tell us which tax law applies to this company.

To expose this dependency, stipulate that the parent's margin of 10 is taxed
now at one quarter, that the subsidiary keeps a tax base of 30, and that the
same rate applies when it uses the deduction. The parent therefore has a
current tax expense and external payable of $5/2$. The group inventory has
carrying amount 20 and tax base 30, leaving a deductible temporary difference
of 10. If the supplied recoverability conditions permit recognition, the
deferred tax asset is also $5/2$. Its income offsets the current expense in
this example; the external payable remains. If recognition is explicitly
withheld, the asset is zero and group profit is lower by $5/2$. The reference
retains these alternatives separately. If the tax treatment is unknown, it
retains the pretax worksheet and requires the missing tax method before
selecting a final report.

These are supplied tax and attribution premises. Recognition requires the
appropriate taxable-profit assessment, and measurement depends on enacted or
substantively enacted rates and expected recovery
\citep[paras.~24--29, 46--51 and 58]{aasb112dated}. In this downstream example
the parent's current tax and the consolidation tax adjustment are attributed
to parent owners. That premise does not resolve upstream transactions or
every group tax arrangement. The retained source editions support the
conditional design; they do not select an actual entity's reporting package
or tax law.

With the deferred asset recognised, the unsold-inventory case closes with
assets $1629/4$, liabilities 27 and equity $1521/4$. Parent equity is $2749/8$
and NCI is $293/8$. Group cash remains 270, whether the internal payable is
settled or outstanding. Inventory and the extra opening parent capital account
for 20 of the change from the acquisition-only case; the separate tax asset
and payable account for $5/2$ on each side. No balancing amount is needed.

\subsection{An outside sale releases only the corresponding part of the margin}

Now let the subsidiary sell half the lot to an outside customer for 20 cash
before the year closes. The example stipulates a qualifying customer contract,
transfer of control and a fixed price, with no returns or additional obligations.
A cash receipt alone would not establish those facts
\citep[paras.~9, 31--38 and 47]{aasb15dated}. The subsidiary recognises revenue
20 and inventory cost 15. Inventory cost is expensed when the related revenue
is recognised \citep[para.~34]{aasb102compiled2021}.

The group still eliminates the full internal sale of 30. Only half of its
margin remains in unsold inventory, so the closing unrealised-profit
adjustment is 5. The remaining group inventory is 10, while the sold half
costs the group 10. The outside sale thus adds group pretax profit of 10:
revenue 20 less the original group cost 10. Of this increment, the
subsidiary's own sale contributes 5 and the release of the parent's margin
contributes 5.

Continue the supplied tax premises. The parent's current tax remains $5/2$;
the subsidiary's profit of 5 adds current tax $5/4$. Only the unsold margin
supports the remaining deferred asset, now $5/4$. Relative to the acquisition
case without inventory trading, the group earns an additional after-tax
profit of $15/2$. Half arises in the subsidiary after its current tax, and
half belongs to parent owners after the downstream consolidation adjustments.
Ordinary NCI therefore participates in the subsidiary's $15/4$ increment,
not in the parent's released margin.

The full closing comparison gives group profit 39: owners receive $483/16$
and NCI receives $141/16$. Assets are 416, liabilities are $113/4$, and equity
is $1551/4$. Parent equity of $5603/16$ plus NCI of $601/16$ gives the same
total. Cash is 290, including the outside receipt of 20. That operating receipt
remains in the group cash report even if the separate internal payment of 30
is also made.

\subsection{The dated journals make the closing report explainable}

The comparison retains every legal entry, including the supplied opening
trial balances, the explicitly stipulated pre-control receipt and interim
close, acquisition consideration, controlled-period receipts and depreciation,
internal sale and settlement, outside sale and current taxes. Replaying these
dated journals recovers the opening, acquisition and closing accounts. Their
origins distinguish an outside receipt from a reciprocal internal payment;
the consolidation entries remain in a separate worksheet. A set of closing
balances could not establish that history on its own.

No resale, half resale and full resale are compared under each selected NCI
election and consideration mix. Unknown tax preserves the missing dependency
instead of completing the report with zero. These results remain a one-lot,
same-period downstream example with supplied revenue, tax and ordinary-rights
premises.\footnote{The complete downstream comparison is retained in
\path{data/accounting-master/acquisition-intragroup-20261027/}; the dated
journal and resale comparison is in
\path{data/accounting-master/acquisition-intragroup-realisation-20261027/}.
Both preserve full independent account, journal and report expectations and
keep source, posting and report authorisation false.}
Carrying unsold inventory into another year creates a new question: the next
report must release profit from opening group equity and carry its tax basis,
rather than repeat the original acquisition or eliminate a sale that occurred
in the prior year. That requires a separate two-year comparison.
